代码里写着 \texttt{x = inv(A) @ b}。它能跑，结果看着也对。

换一个稍微病态一点的 \(A\)，再和 \texttt{solve(A, b)} 比一下，差别就出来了：解方程比先求逆再相乘快一倍，而且误差界里少了一个条件数因子。两行代码在数学上是同一个 \(x\)，在浮点里不是同一件事——这个观察在上一章已经出现过，这里要把它系统化。

线性代数课上，\(A^{-1}b\)、LU、QR、SVD 都在描述同一个解，所以选哪个是审美问题。到了计算机上它们分道扬镳：一个 1000 阶的对称正定矩阵，带主元 LU 要约 \(3.3\times10^8\) 次运算，Cholesky 只要一半，而且不必选主元；若矩阵还是带状的，两者都能降到 \(O(nw^2)\)。\textbf{矩阵的结构不是背景信息，它就是可以兑现的计算资源。}

于是这一章要回答的问题很具体：手上这个矩阵有什么结构、规模有多大、你真正想要的是什么，三者合起来决定该用哪种分解。

\begin{center}
\begin{tikzpicture}[x=1cm,y=1cm,font=\small,>=stealth]
  \draw[black!45,fill=blue!9] (0,0.2) rectangle (2.9,1.6);
  \node[align=center] at (1.45,0.9) {手上的\\矩阵};
  \draw[black!45,fill=green!8] (4.4,3.35) rectangle (12.6,4.25);
  \node[align=left,anchor=west] at (4.6,3.8) {一般方阵，反复解 \(\;\to\;\) 带主元 LU（第一节）};
  \draw[black!45,fill=green!8] (4.4,2.2) rectangle (12.6,3.1);
  \node[align=left,anchor=west] at (4.6,2.65) {长方形，最小二乘 \(\;\to\;\) QR（第二节）};
  \draw[black!45,fill=green!8] (4.4,1.05) rectangle (12.6,1.95);
  \node[align=left,anchor=west] at (4.6,1.5) {对称正定 \(\;\to\;\) Cholesky（第三节）};
  \draw[black!45,fill=orange!9] (4.4,-0.1) rectangle (12.6,0.8);
  \node[align=left,anchor=west] at (4.6,0.35) {太大装不下，只能做矩阵乘向量 \(\;\to\;\) Krylov（第四节）};
  \draw[black!45,fill=orange!9] (4.4,-1.25) rectangle (12.6,-0.35);
  \node[align=left,anchor=west] at (4.6,-0.8) {想要谱或主方向 \(\;\to\;\) 特征值算法、SVD（第五、六节）};
  \draw[black!45,fill=orange!9] (4.4,-2.4) rectangle (12.6,-1.5);
  \node[align=left,anchor=west] at (4.6,-1.95) {一次只能读一行 \(\;\to\;\) Kaczmarz（第七节）};
  \foreach \y in {3.8,2.65,1.5,0.35,-0.8,-1.95}
    \draw[->,black!50] (2.95,0.9) -- (4.35,\y);
  \node[align=center,font=\footnotesize,black!65] at (1.45,-1.3) {绿色：整个矩阵\\装得进内存\\[4pt]橙色：装不进，\\或不必装进};
\end{tikzpicture}

\emph{七节不是七种并列的技巧，而是同一个问题在不同结构与不同规模下的回答。分岔点只有三个：矩阵长什么样、大到什么程度、你要的是解还是谱。}
\end{center}

这张分岔图管的是前七节。最后三节不在图上，因为它们不是第七条分支：它们不换分解方式，换的是结论的种类。前七节的每一条结论都形如``条件满足就一定成立''，而随机草图、随机化 SVD 给出的是``以概率至少 \(1-\delta\) 成立''。规模再大一个量级时，能拿到的往往只剩下后一种，所以要先看清这两句话差在哪里。

\subsection{四条判断，读完这一章要能自己用}

\textbf{分解一次，复用多次。}\(O(n^3)\) 的代价花在分解上，之后每个新的右端项只要 \(O(n^2)\)。这就是为什么标准库把分解与求解拆成两个函数，也是为什么``求逆''这个动作在实践中几乎从不出现——你要的从来不是 \(A^{-1}\)，是 \(x\)。

\textbf{结构决定代价，也决定可信度。}对称正定不只省一半运算，它还免掉选主元、并让分解失败成为一个有用的信号。稀疏不只省内存，它把直接法整个换成了迭代法，而迭代法的收敛速度不看条件数看谱的分布。每次利用结构，都是在用一个假设换取运算量的下降——所以也要问一句：这个假设万一不成立，会以什么方式暴露出来。

\textbf{残差小不等于解准。}方程几乎被满足，与未知量几乎正确，中间隔着一个条件数。上一章那条``最终误差 \(\approx\) 后向误差 \(\times\) 条件数''的关系在这里反复出现：QR 之所以优于正规方程，正是因为后者把条件数平方了。检查工作时不要只看残差。

\textbf{保证分强弱，用之前先认清手上是哪一种。}同一个残差数字，可以来自一条对一切输入都成立的确定性界，也可以来自一条只在 \(1-\delta\) 的运行里成立的概率界。两者都能写进报告，但能支撑的动作不一样：后者不足以支撑一次需要复现的审计，也不足以支撑一个最坏情况的承诺。规模逼得你换到概率保证时，这件事要说出来，而不是让它藏在实现细节里。

\subsection{怎么读这十节}

前三节是主线，建立最常用的对应关系，建议连着读完。\hyperref[sec:08-Numerical-Analysis/05-Numerical-Linear-Algebra/01]{``线性方程为何分解而不求逆''}讲分解与选主元的机制，\hyperref[sec:08-Numerical-Analysis/05-Numerical-Linear-Algebra/02]{``QR 与最小二乘的稳定解法''}给出回归问题该怎么解，\hyperref[sec:08-Numerical-Analysis/05-Numerical-Linear-Algebra/03]{``Cholesky：利用正定结构''}说明一个结构假设能买到多少。

\hyperref[sec:08-Numerical-Analysis/05-Numerical-Linear-Algebra/06]{``SVD、低秩与数值秩''}可以提前读——它解释了``秩''在浮点世界里为什么不是一个干脆的整数，而这件事在做降维和低秩近似时天天遇到。

后面三节按你遇到的规模瓶颈取用：矩阵大到无法完整分解但能做矩阵乘向量，读\hyperref[sec:08-Numerical-Analysis/05-Numerical-Linear-Algebra/04]{``大型稀疏系统与 Krylov 迭代''}；需要的是特征值而不是解，读\hyperref[sec:08-Numerical-Analysis/05-Numerical-Linear-Algebra/05]{``从幂迭代到 QR 特征值算法''}；连一整个矩阵都放不进内存、一次只能读一行，读\hyperref[sec:08-Numerical-Analysis/05-Numerical-Linear-Algebra/07]{``Kaczmarz：逐行投影大型线性系统''}——它与随机梯度下降是同一种结构，每次只看一个样本。

最后三节自成一组，讲的是大规模场景里真正在用的那一套：矩阵大到连一次完整分解都排不上日程时，把随机性主动引进算法。\hyperref[sec:08-Numerical-Analysis/05-Numerical-Linear-Algebra/08]{``用随机草图把大矩阵压小''}用一个随机矩阵从左边压掉行数，把上千万行的最小二乘换成几万行的；\hyperref[sec:08-Numerical-Analysis/05-Numerical-Linear-Algebra/09]{``随机化 SVD：用随机向量探到主子空间''}把随机矩阵挪到右边，两遍数据就取出前 \(k\) 个主方向，正好接上第六节那句``只需要前 \(k\) 个方向时完整 SVD 是浪费''；\hyperref[sec:08-Numerical-Analysis/05-Numerical-Linear-Algebra/10]{``随机化给的是概率保证''}把前两节共有的那半句话单独拿出来说清楚——什么时候一个 \(1-\delta\) 的承诺够用，什么时候它根本不该被当成承诺。这三节读起来不需要前七节的全部细节，但第二节的 QR 和第六节的 SVD 是它们的落脚点，先读那两节更顺。

上一章是\hyperref[ch:064]{``数值微分与积分：同一份采样，为何难度相反''}，那里的主角是步长；这里的主角换成了矩阵结构，但追问方式没变：这个做法在什么条件下可信，条件塌掉时会以什么形式出错。
